% Options for packages loaded elsewhere
\PassOptionsToPackage{unicode}{hyperref}
\PassOptionsToPackage{hyphens}{url}
\documentclass[
  12pt,
]{ctexart}
\usepackage{xcolor}
\usepackage{amsmath,amssymb}
\setcounter{secnumdepth}{-\maxdimen} % remove section numbering
\usepackage{iftex}
\ifPDFTeX
  \usepackage[T1]{fontenc}
  \usepackage[utf8]{inputenc}
  \usepackage{textcomp} % provide euro and other symbols
\else % if luatex or xetex
  \usepackage{unicode-math} % this also loads fontspec
  \defaultfontfeatures{Scale=MatchLowercase}
  \defaultfontfeatures[\rmfamily]{Ligatures=TeX,Scale=1}
\fi
\usepackage{lmodern}
\ifPDFTeX\else
  % xetex/luatex font selection
\fi
% Use upquote if available, for straight quotes in verbatim environments
\IfFileExists{upquote.sty}{\usepackage{upquote}}{}
\IfFileExists{microtype.sty}{% use microtype if available
  \usepackage[]{microtype}
  \UseMicrotypeSet[protrusion]{basicmath} % disable protrusion for tt fonts
}{}
\makeatletter
\@ifundefined{KOMAClassName}{% if non-KOMA class
  \IfFileExists{parskip.sty}{%
    \usepackage{parskip}
  }{% else
    \setlength{\parindent}{0pt}
    \setlength{\parskip}{6pt plus 2pt minus 1pt}}
}{% if KOMA class
  \KOMAoptions{parskip=half}}
\makeatother
\usepackage{longtable,booktabs,array}
\usepackage{calc} % for calculating minipage widths
% Correct order of tables after \paragraph or \subparagraph
\usepackage{etoolbox}
\makeatletter
\patchcmd\longtable{\par}{\if@noskipsec\mbox{}\fi\par}{}{}
\makeatother
% Allow footnotes in longtable head/foot
\IfFileExists{footnotehyper.sty}{\usepackage{footnotehyper}}{\usepackage{footnote}}
\makesavenoteenv{longtable}
\setlength{\emergencystretch}{3em} % prevent overfull lines
\providecommand{\tightlist}{%
  \setlength{\itemsep}{0pt}\setlength{\parskip}{0pt}}
\usepackage{bookmark}
\IfFileExists{xurl.sty}{\usepackage{xurl}}{} % add URL line breaks if available
\urlstyle{same}
\hypersetup{
  hidelinks,
  pdfcreator={LaTeX via pandoc}}

\author{SCX}
\date{2026}

\begin{document}

\section{SCX Self-Evolution as a Discrete Dynamical
System}\label{scx-self-evolution-as-a-discrete-dynamical-system}

\begin{quote}
\textbf{Version}: 2026-06-28 \textbar{} \textbf{Status}: Theoretical
development \textbar{} \textbf{Audit}: Pre-verification
\textbf{Purpose}: Formalize the SCX self-evolution loop as a discrete
dynamical system, study fixed points, Lyapunov stability, attractor
structure, and phase portraits. \textbf{Prerequisites}: Document 01
(Symbol System and Problem Setup)
\end{quote}

\begin{center}\rule{0.5\linewidth}{0.5pt}\end{center}

\subsection{Table of Contents}\label{table-of-contents}

\begin{enumerate}
\def\labelenumi{\arabic{enumi}.}
\tightlist
\item
  \hyperref[1-discrete-dynamical-system-formulation]{Discrete Dynamical
  System Formulation}
\item
  \hyperref[2-state-space-topology-and-metric-on-s_t]{State Space
  Topology and Metric on S\_t}
\item
  \hyperref[3-orbits-under-phi]{Orbits Under Phi}
\item
  \hyperref[4-attractors-and-omega-limit-sets]{Attractors and
  Omega-Limit Sets}
\item
  \hyperref[5-fixed-points]{Fixed Points}
\item
  \hyperref[6-existence-conditions-for-fixed-points]{Existence
  Conditions for Fixed Points}
\item
  \hyperref[7-lyapunov-function-candidate]{Lyapunov Function Candidate}
\item
  \hyperref[8-monotonicity-analysis]{Monotonicity Analysis}
\item
  \hyperref[9-phase-portrait]{Phase Portrait}
\item
  \hyperref[10-connection-to-existing-theorems-1-3]{Connection to
  Existing Theorems 1-3}
\item
  \hyperref[11-summary-of-proven-vs-conjectured-claims]{Summary of
  Proven vs.~Conjectured Claims}
\end{enumerate}

\begin{center}\rule{0.5\linewidth}{0.5pt}\end{center}

\subsection{1. Discrete Dynamical System
Formulation}\label{discrete-dynamical-system-formulation}

\subsubsection{1.1 Basic Formulation}\label{basic-formulation}

Let \(z_t = (S_t, M_t, \theta_t) \in \mathcal{Z}\) be the evolution
state at time \(t\) as defined in Document 01, Section 5. The
self-evolution loop is a discrete dynamical system:

\[z_{t+1} = \Phi(z_t), \qquad t = 0, 1, 2, \dots\]

with initial condition \(z_0 = (S_0, M_0, \theta_0)\) determined by the
static SCX initialization.

\textbf{Definition 16 (Discrete-Time Dynamical System).} The tuple
\((\mathcal{Z}, \Phi, \mathbb{N}_0)\) constitutes a discrete-time
dynamical system where: -
\(\mathcal{Z} = \mathcal{F} \times \mathcal{M} \times \Theta\) is the
state space (Section 5 of Document 01) -
\(\Phi: \mathcal{Z} \to \mathcal{Z}\) is the update operator (Section 9
of Document 01) - \(\mathbb{N}_0 = \{0, 1, 2, \dots\}\) is the discrete
time index

\subsubsection{1.2 Decomposition into Component
Maps}\label{decomposition-into-component-maps}

The operator \(\Phi\) decomposes into three coupled maps:

\[\Phi(S, M, \theta) = \bigl( \Phi_S(S, M, \theta),\; \Phi_M(S),\; \Phi_\theta(S, M) \bigr)\]

where the coupling structure is:

\begin{verbatim}
S_t ──Φ_M──→ M_{t+1}
S_t ──Φ_θ──→ θ_{t+1}  (via M_{t+1})
S_t, M_{t+1}, θ_{t+1} ──Φ_S──→ S_{t+1}
\end{verbatim}

This coupling is \textbf{feed-forward}: \(S_t\) influences \(M_{t+1}\)
and \(\theta_{t+1}\), which in turn influence \(S_{t+1}\). There is no
instantaneous feedback loop.

\textbf{Proposition 3 (Causal Structure).} The system's information flow
is a directed acyclic graph within each timestep:
\(S_t \to M_{t+1} \to \theta_{t+1}\), and
\((S_t, M_{t+1}, \theta_{t+1}) \to S_{t+1}\). Consequently, the Jacobian
of \(\Phi\) is block-triangular.

\emph{Proof.} By inspection of the update definitions: \(\Phi_M\)
depends only on \(S_t\) (not on \(M_t\) or \(\theta_t\));
\(\Phi_\theta\) depends on \(M_{t+1}\) and implicitly on \(S_t\) through
\(M_{t+1}\); \(\Phi_S\) depends on all three. \(\square\)

\begin{center}\rule{0.5\linewidth}{0.5pt}\end{center}

\subsection{\texorpdfstring{2. State Space Topology and Metric on
\(S_t\)}{2. State Space Topology and Metric on S\_t}}\label{state-space-topology-and-metric-on-s_t}

\subsubsection{2.1 Metric Structure}\label{metric-structure}

Recall from Document 01 the metric on \(\mathcal{F}\):

\[d_{\mathcal{F}}(S, S') = \mathbb{E}_{(x,y) \sim \mathcal{D}}\bigl[ |S(x,y) - S'(x,y)| \bigr]\]

and the product metric on \(\mathcal{Z}\):

\[d_{\mathcal{Z}}(z, z') = d_{\mathcal{F}}(S, S') + d_{\mathcal{M}}(M, M') + \|\theta - \theta'\|_2\]

\textbf{Proposition 4 (Completeness).} If \(\mathcal{D}\) has full
support on \(\mathcal{X}_0 \times \mathcal{Y}\), then
\((\mathcal{F}, d_{\mathcal{F}})\) is a complete metric space. The
product space \((\mathcal{Z}, d_{\mathcal{Z}})\) is complete.

\emph{Proof sketch.} \(L^1(\mathcal{D})\) is complete. The finite set
\(\mathcal{M}\) with the discrete topology is trivially complete.
\(\Theta \subseteq \mathbb{R}^{d_\theta}\) is complete if closed. The
product of complete metric spaces is complete. \(\square\)

\subsubsection{2.2 Compactness
Considerations}\label{compactness-considerations}

For existence of fixed points and attractors, compactness of the state
space is desirable.

\textbf{Proposition 5 (Compactness of \(\mathcal{F}\) under
Regularization).} If we restrict \(\mathcal{F}\) to the set of
functions: - bounded in \(L^\infty\) norm (already \([0,1]\)-valued) -
with uniformly bounded Lipschitz constant \(L_S < \infty\) with respect
to input distance then \(\mathcal{F}\) is compact in \(d_{\mathcal{F}}\)
by the Arzela-Ascoli theorem (on compact subsets of \(\mathcal{X}_0\)).

\emph{Practical implication.} In implementation, neural network
gatekeepers satisfy Lipschitz bounds (e.g., via spectral normalization).
Theoretical analysis often assumes such regularization to ensure
well-posedness.

\begin{center}\rule{0.5\linewidth}{0.5pt}\end{center}

\subsection{\texorpdfstring{3. Orbits Under
\(\Phi\)}{3. Orbits Under \textbackslash Phi}}\label{orbits-under-phi}

\subsubsection{3.1 Definition}\label{definition}

\textbf{Definition 17 (Orbit).} The forward orbit of an initial state
\(z_0 \in \mathcal{Z}\) under \(\Phi\) is the sequence:

\[\mathcal{O}^+(z_0) = \{ z_0, \Phi(z_0), \Phi^{(2)}(z_0), \Phi^{(3)}(z_0), \dots \}\]

The backward orbit (if \(\Phi\) is invertible) is defined analogously
but is not guaranteed.

\subsubsection{3.2 Types of Orbits}\label{types-of-orbits}

Based on the system's behavior, we classify orbits into three types:

\textbf{Definition 18 (Orbit Classification).}

\begin{enumerate}
\def\labelenumi{\arabic{enumi}.}
\tightlist
\item
  \textbf{Convergent orbit}:
  \(\lim_{t\to\infty} d_{\mathcal{Z}}(z_t, z^*) = 0\) for some
  \(z^* \in \mathcal{Z}\).
\item
  \textbf{Periodic orbit}: \(\exists T > 0\) such that \(z_{t+T} = z_t\)
  for all sufficiently large \(t\), with no smaller period.
\item
  \textbf{Chaotic orbit}: The orbit is bounded but does not converge and
  is not periodic; sensitive dependence on initial conditions may occur.
\end{enumerate}

\textbf{Conjecture 2 (Orbit Classification for SCX).} Under Assumptions
B1-B3, all orbits of the SCX self-evolution system are convergent. No
periodic or chaotic orbits exist.

\emph{Status: \textbf{Conjecture}. The monotonicity of \(M_t\) (B1)
prevents cycling through memory loss, but stochastic gradient updates in
\(\Phi_S\) and \(\Phi_\theta\) could, in principle, produce limit
cycles. The conjecture asserts that the Lyapunov function (Section 7)
eliminates this possibility.}

\subsubsection{3.3 Properties of Orbits}\label{properties-of-orbits}

\textbf{Proposition 6 (Boundedness of Orbits).} Under Assumption B2
(bounded memory growth), any orbit \(\mathcal{O}^+(z_0)\) is bounded in
\(\mathcal{Z}\).

\emph{Proof.} \(S_t\) is bounded in \([0,1]\) by construction,
\(\|S_t\|_\infty \leq 1\). \(\theta_t\) is bounded if the loss landscape
has bounded sublevel sets (e.g., under L2 regularization). \(M_t\) has
bounded size by B2. Hence \(d_{\mathcal{Z}}(z_0, z_t)\) is bounded for
all \(t\). \(\square\)

\begin{center}\rule{0.5\linewidth}{0.5pt}\end{center}

\subsection{\texorpdfstring{4. Attractors and \(\omega\)-Limit
Sets}{4. Attractors and \textbackslash omega-Limit Sets}}\label{attractors-and-omega-limit-sets}

\subsubsection{4.1 Definition}\label{definition-1}

\textbf{Definition 19 (\(\omega\)-Limit Set).} For a trajectory
\(z_t = \Phi^{(t)}(z_0)\), the \(\omega\)-limit set is:

\[\omega(z_0) = \bigl\{ z \in \mathcal{Z} : \exists \{t_k\}_{k=1}^\infty \text{ with } t_k \to \infty \text{ and } \lim_{k\to\infty} d_{\mathcal{Z}}(z_{t_k}, z) = 0 \bigr\}\]

The \(\alpha\)-limit set (for backward orbits) is defined analogously.

\textbf{Definition 20 (Attractor).} A set \(A \subseteq \mathcal{Z}\) is
an attractor if: 1. \(A\) is invariant: \(\Phi(A) = A\). 2. There exists
a neighborhood \(U\) of \(A\) such that \(\omega(z) \subseteq A\) for
all \(z \in U\). 3. \(A\) is minimal with respect to properties 1 and 2.

The \textbf{basin of attraction} of \(A\) is:

\[\mathcal{B}(A) = \{ z \in \mathcal{Z} : \omega(z) \subseteq A \}\]

\subsubsection{4.2 Candidate Attractors for
SCX}\label{candidate-attractors-for-scx}

\textbf{Proposition 7 (Fixed Points as Attractors).} If \(z^*\) is an
asymptotically stable fixed point (Definition 21), then \(\{z^*\}\) is
an attractor with basin \(\mathcal{B}(z^*)\) containing an open
neighborhood of \(z^*\).

The primary candidate attractor for the SCX self-evolution system is the
\textbf{perfect gatekeeper}:

\[S^*(x,y) = \mathbf{1}\{y = f^*(x)\}\]

i.e., the gatekeeper that perfectly identifies whether the observed
label matches the true oracle. However, achieving this fixed point
requires: 1. The NEP student converges to \(f^*\) (Assumption B3) 2. The
memory bank contains sufficient clean samples

\textbf{Conjecture 3 (Unique Attractor).} Under Assumptions B1-B3 and
the uniform noise assumption (A4), the SCX self-evolution system has a
unique attractor \(z^* = (S^*, M^*, \theta^*)\) where
\(S^*(x,y) = \mathbf{1}\{y = f^*(x)\}\), \(M^*\) is the infinite memory
bank containing all sufficiently clean samples, and \(\theta^*\)
parameterizes \(f^*\).

\emph{Status: \textbf{Conjecture}. The uniqueness depends on the loss
landscape of the gatekeeper and NEP being convex, which is not
guaranteed for neural network parameterizations. However, if the
function spaces are sufficiently expressive and the losses are proper
scoring rules, the fixed point is unique in function space.}

\begin{center}\rule{0.5\linewidth}{0.5pt}\end{center}

\subsection{5. Fixed Points}\label{fixed-points}

\subsubsection{5.1 Definition}\label{definition-2}

\textbf{Definition 21 (Fixed Point).} A state
\(z^* = (S^*, M^*, \theta^*) \in \mathcal{Z}\) is a fixed point of the
dynamics if:

\[\Phi(z^*) = z^*\]

Equivalently:

\[\Phi_S(S^*, M^*, \theta^*) = S^*\] \[\Phi_M(S^*, M^*) = M^*\]
\[\Phi_\theta(S^*, M^*) = \theta^*\]

\textbf{Definition 22 (Stability of Fixed Points).} A fixed point
\(z^*\) is: - \textbf{Lyapunov stable} if \(\forall \varepsilon > 0\),
\(\exists \delta > 0\) such that \(d_{\mathcal{Z}}(z_0, z^*) < \delta\)
implies \(d_{\mathcal{Z}}(z_t, z^*) < \varepsilon\) for all
\(t \geq 0\). - \textbf{Asymptotically stable} if it is Lyapunov stable
and \(\exists \delta > 0\) such that
\(d_{\mathcal{Z}}(z_0, z^*) < \delta\) implies
\(\lim_{t\to\infty} d_{\mathcal{Z}}(z_t, z^*) = 0\). - \textbf{Unstable}
if it is not Lyapunov stable.

\subsubsection{5.2 Fixed Point Equations}\label{fixed-point-equations}

\textbf{Fixed point equation for \(S^*\):}

\[S^*(x, y) = \text{Update}_{\text{gate}}\bigl(S^*, M^*, f_{\theta^*}\bigr)(x, y)\]

For a gradient-based update
\(\text{Update}_{\text{gate}}(S, M, f_\theta) = S - \alpha \nabla_S \hat{L}_{\text{gate}}(S; M)\),
this implies:

\[\nabla_S \hat{L}_{\text{gate}}(S^*; M^*) = 0\]

i.e., \(S^*\) is a stationary point of the gatekeeper loss on the memory
bank \(M^*\).

\textbf{Fixed point equation for \(M^*\):}

\[M^* = M^* \cup \{(x, y, S^*(x,y), S^*(x,y)) : (x,y) \in q, \; S^*(x,y) > \delta_{\text{store}}\}\]

This implies:

\[\{ (x,y) \in \text{supp}(\mathcal{D}) : S^*(x,y) > \delta_{\text{store}} \} \subseteq M^*\]

i.e., \(M^*\) contains all samples that the fixed-point gatekeeper
judges reliable. Since memory is monotone (B1), \(M^*\) is the limit of
the growing sequence.

\textbf{Fixed point equation for \(\theta^*\):}

\[\theta^* = \arg\min_{\theta \in \Theta} \frac{1}{|M^*|} \sum_{(x_i, y_i, \cdot, \cdot) \in M^*} \ell_{\text{nep}}(f_\theta(x_i), y_i)\]

i.e., \(\theta^*\) minimizes the NEP loss on \(M^*\).

\subsubsection{5.3 Trivial Fixed Points}\label{trivial-fixed-points}

\textbf{Proposition 8 (Trivial Fixed Points).} The following are fixed
points of \(\Phi\):

\begin{enumerate}
\def\labelenumi{\arabic{enumi}.}
\item
  \textbf{Zero gatekeeper}: \(S_0 \equiv 0\), \(M_0 = \varnothing\),
  \(\theta_0\) arbitrary. The system never stores any samples and never
  updates. (Unstable --- any positive \(S\) triggers storage.)
\item
  \textbf{Perfect gatekeeper}: \(S^*(x, y) = \mathbf{1}\{y = f^*(x)\}\),
  \(M^* = \{(x, f^*(x), 1, 1) : x \in \mathcal{X}_0\}\),
  \(f_{\theta^*} = f^*\). All samples are judged correctly; nothing
  changes. (Asymptotically stable if the NEP converges.)
\item
  \textbf{Constant gatekeeper}: \(S(x, y) \equiv c\) for some
  \(c \neq \delta_{\text{store}}\). If \(c < \delta_{\text{store}}\), no
  samples are stored (like 1). If \(c > \delta_{\text{store}}\), all
  samples are stored, \(M\) grows to the entire dataset.
\end{enumerate}

\emph{Proof sketch.} Each satisfies the fixed point equations by
inspection. For stability claims, see Section 8. \(\square\)

\begin{center}\rule{0.5\linewidth}{0.5pt}\end{center}

\subsection{6. Existence Conditions for Fixed
Points}\label{existence-conditions-for-fixed-points}

\subsubsection{6.1 General Existence}\label{general-existence}

\textbf{Theorem 4 (Existence of Fixed Points).} Suppose: 1. \(\Phi\) is
continuous on \(\mathcal{Z}\) in the product metric \(d_{\mathcal{Z}}\).
2. \(\mathcal{Z}\) is compact in \(d_{\mathcal{Z}}\). 3. \(\Phi\) maps
\(\mathcal{Z}\) into itself.

Then \(\Phi\) has at least one fixed point in \(\mathcal{Z}\).

\emph{Proof.} Apply the Brouwer fixed-point theorem (for compact convex
subsets of a Banach space) or Schauder-Tychonoff (for locally convex
topological vector spaces). The product structure
\(\mathcal{Z} = \mathcal{F} \times \mathcal{M} \times \Theta\) with
\(\mathcal{F}\) compact (Proposition 5) and \(\Theta\) compact (under
boundedness), and \(\mathcal{M}\) finite, satisfies the hypotheses.
\(\square\)

\textbf{Remark.} Condition 1 (continuity of \(\Phi\)) is non-trivial:
the \(\arg\min\) in \(\Phi_\theta\) is set-valued if the minimum is not
unique. Formalizing \(\Phi\) as a selection from a correspondence (upper
hemicontinuous) allows using the Kakutani fixed-point theorem instead.

\subsubsection{6.2 Specific Existence for
SCX}\label{specific-existence-for-scx}

\textbf{Theorem 5 (Fixed Point Existence for SCX with Convex Loss).}
Suppose: 1. The gatekeeper loss \(L_{\text{gate}}(S; M)\) is strictly
convex in \(S\) for any \(M\). 2. The NEP loss
\(L_{\text{nep}}(\theta; M)\) is strictly convex in \(\theta\) for any
\(M\). 3. The memory bank growth saturates:
\(\lim_{t\to\infty} \mathbb{P}(S_t(x,y) > \delta_{\text{store}}) = 0\)
or \(1\) for each \((x,y)\).

Then there exists a unique fixed point \(z^* = (S^*, M^*, \theta^*)\).

\emph{Proof sketch.} Convexity ensures unique minimizers for \(S^*\) and
\(\theta^*\), making \(\Phi\) a function (not correspondence). Under the
saturation condition, \(M^*\) is the limit of the monotone sequence. The
contraction mapping principle (if \(\Phi\) is a contraction) or Schauder
fixed point gives existence. \(\square\)

\emph{Status: \textbf{Rigorous} under the stated convexity assumptions.
In practice, neural network losses are non-convex, so the conditions are
violated. The theorem provides an idealized benchmark.}

\subsubsection{6.3 Necessary Condition for Non-Trivial Fixed
Point}\label{necessary-condition-for-non-trivial-fixed-point}

\textbf{Proposition 9 (Necessary Condition).} If \(z^*\) is a fixed
point with \(M^* \neq \varnothing\), then the NEP student
\(f_{\theta^*}\) must satisfy:

\[\mathbb{E}_{(x,y) \sim \mathcal{D}}\bigl[ \ell_{\text{nep}}(f_{\theta^*}(x), y) \mid S^*(x,y) > \delta_{\text{store}} \bigr] \leq \mathbb{E}_{(x,y) \sim \mathcal{D}}\bigl[ \ell_{\text{nep}}(f_{\theta^*}(x), y) \mid S^*(x,y) \leq \delta_{\text{store}} \bigr]\]

i.e., the NEP performs no worse on samples the gatekeeper admits than on
samples it rejects.

\emph{Proof.} At a fixed point, the gatekeeper's storage policy and the
NEP's parameters are consistent. If the inequality were reversed, the
gatekeeper would be systematically mislabeling high-loss samples as
reliable, contradicting the stationarity of \(S^*\) (since \(S^*\)
minimizes \(L_{\text{gate}}\) on \(M^*\), and \(M^*\) contains only
admitted samples). \(\square\)

\begin{center}\rule{0.5\linewidth}{0.5pt}\end{center}

\subsection{7. Lyapunov Function --- Explicit Candidate (DEFECT-03/04
Fix)}\label{lyapunov-function-explicit-candidate-defect-0304-fix}

\begin{quote}
\textbf{Status change (2026-06-28).} The original Section 7 defined a
vague Lyapunov candidate
\(V(z_t) = \mathbb{E}[L_{\text{gate}}] + \lambda \cdot \mathbb{E}[L_{\text{nep}}]\)
evaluated on the acceptance-biased distribution. This was found to be
\textbf{underspecified} (DEFECT-03) and to contain a \textbf{selection
bias confound} (DEFECT-04): a decrease in \(V\) could reflect
distribution shift rather than genuine model improvement. The following
explicit candidate on a \textbf{fixed reference set} is adopted as a
concrete target for future proof. The convergence claim is downgraded
from ``Proven'' to \textbf{Conjecture}.
\end{quote}

\subsubsection{7.1 Concrete Lyapunov Function
Candidate}\label{concrete-lyapunov-function-candidate}

\textbf{Definition 23 (Concrete Lyapunov Function Candidate).} Define
\(\Psi: \mathcal{F} \times \Theta \to \mathbb{R}_{\geq 0}\):

\[\boxed{\;\Psi(S_t, \theta_t) = \frac{1}{|M_0|} \sum_{x \in M_0} \bigl(S_t(x, y(x)) - \hat{C}(x)\bigr)^2 \;+\; \lambda \cdot \frac{1}{|V_0|} \sum_{(x,y) \in V_0} \ell(f_{\theta_t}(x), y)\;},\]

where: - \(M_0 \subset \mathcal{X} \times \mathcal{Y}\) is a
\textbf{fixed reference set} (e.g., initial samples before
self-evolution begins, or an external validation set). \(M_0\) does
\textbf{not} grow with \(t\). - \(V_0 \subseteq M_0\) is the subset of
\(M_0\) with externally verified clean labels, -
\(\hat{C}(x) = \frac{1}{M}\sum_{m=1}^M \mathbf{1}\{\ell(f_m(x), y) > \tau\}\)
is the consensus score computed from the \(M\) expert models (fixed), -
\(\lambda > 0\) is a balancing hyperparameter.

\begin{quote}
\textbf{Notation note (B3 fix):} The Lyapunov function is denoted
\(\Psi\) (Psi), NOT \(\Phi\) (Phi). The symbol \(\Phi\) is reserved for
the \textbf{Update Operator} \(\Phi: \mathcal{Z} \to \mathcal{Z}\)
defined in Document 01, Section 9. This resolves the \(\Phi\)-\(\Psi\)
notation collision identified in the cross-file consistency audit.
\end{quote}

\textbf{Why a fixed reference set?} The original Lyapunov candidate
evaluated both gatekeeper and student losses on the acceptance-biased
distribution \(P_{S_t}\). This created a \textbf{selection bias
confound} (DEFECT-04/13): the gatekeeper can decrease its apparent loss
by becoming more selective and admitting only ``easy'' samples, without
improving its true discriminative ability. By evaluating on a fixed
reference set \(M_0\) that does not change with \(t\), we ensure that a
decrease in \(\Phi\) reflects genuine improvement.

\subsubsection{7.2 Component
Decomposition}\label{component-decomposition}

The Lyapunov function decomposes into two interpretable terms:

\textbf{Gatekeeper term (\(\Psi_{\text{gate}}\)).} Measures the squared
discrepancy between the gatekeeper's score \(S_t(x, y(x))\) and the
fixed expert consensus \(\hat{C}(x)\) on the reference set. This term
decreases when the gatekeeper's judgments align better with expert
consensus:
\[\Psi_{\text{gate}}(S_t) = \frac{1}{|M_0|} \sum_{x \in M_0} (S_t(x, y(x)) - \hat{C}(x))^2.\]

\textbf{Student term (\(\Psi_{\text{student}}\)).} Measures the NEP
student's prediction error on the verified-clean subset \(V_0\). This
term decreases when the student learns better representations:
\[\Psi_{\text{student}}(\theta_t) = \frac{1}{|V_0|} \sum_{(x,y) \in V_0} \ell(f_{\theta_t}(x), y).\]

\subsubsection{7.3 Properties}\label{properties}

\textbf{Proposition 10 (Properties of \(\Psi\)).} 1.
\textbf{Non-negativity}: \(\Psi(S, \theta) \geq 0\) for all
\((S, \theta)\), with equality iff \(S_t(x, y(x)) = \hat{C}(x)\) for all
\(x \in M_0\) and \(f_{\theta_t}(x) = y\) for all \((x,y) \in V_0\). 2.
\textbf{Boundedness}: Under Assumption A3 (bounded loss
\(\ell \leq B\)), \(\Psi(S, \theta) \leq 1 + \lambda B < \infty\) since
\(S_t \in [0,1]\) and \(\hat{C} \in [0,1]\). 3. \textbf{Continuity}:
\(\Psi\) is continuous in \((S, \theta)\) under the product metric
\(d_{\mathcal{Z}}\) if \(\ell\) is continuous and \(M_0\) is finite.

\emph{Proof.} (1) Sums of squares and non-negative losses. (2)
\(|S_t - \hat{C}| \leq 1\), so the squared term \(\leq 1\);
\(\ell \leq B\) by A3. (3) Composition of continuous functions; for
finite \(M_0\), the empirical average is continuous in each argument.
\(\square\)

\subsubsection{7.4 Descent Analysis --- Honest
Assessment}\label{descent-analysis-honest-assessment}

\textbf{Definition 24 (Lyapunov Descent Property).} \(\Psi\) is a
Lyapunov function for the dynamical system \((\mathcal{Z}, \Phi)\)
(where \(\Phi\) is the Update Operator) if:

\[\mathbb{E}[\Psi(S_{t+1}, \theta_{t+1}) \mid \mathcal{F}_t] \leq \Psi(S_t, \theta_t) - \eta_t, \quad \eta_t \geq 0,\]

with \(\eta_t > 0\) except at fixed points.

\textbf{Theorem 6 (Lyapunov Descent --- CONDITIONAL, CONJECTURED).}
Under the following jointly sufficient conditions, the candidate
\(\Psi\) satisfies the descent property:

\begin{enumerate}
\def\labelenumi{\arabic{enumi}.}
\tightlist
\item
  \textbf{Two-timescale separation (C6')}: \(\beta_t = o(\alpha_t)\) ---
  gatekeeper updates are asymptotically slower than student updates.
\item
  \textbf{Student gradient descent (C2, C4)}: Standard SGD with
  Robbins-Monro rates on \(\alpha_t\).
\item
  \textbf{Gatekeeper moves toward consensus (C3, C7)}: SCXUpdate reduces
  the discrepancy \(|S_t - \hat{C}|\) when applied to new data.
\item
  \textbf{Random exploration (C8)}: A fraction \(\varepsilon > 0\) of
  accepted samples are chosen uniformly at random regardless of \(S_t\)
  score, ensuring the reference distribution remains represented.
\item
  \textbf{Annealing acceptance threshold (C9)}: \(\gamma_t \to 0.5\) as
  \(t \to \infty\), starting from a permissive initial value
  \(\gamma_0 < 0.5\).
\end{enumerate}

\emph{Proof sketch (gaps acknowledged).} Decompose the one-step change:

\[\Psi(S_{t+1}, \theta_{t+1}) - \Psi(S_t, \theta_t) = \underbrace{[\Psi_{\text{gate}}(S_{t+1}) - \Psi_{\text{gate}}(S_t)]}_{\text{(A)}} + \lambda \underbrace{[\Psi_{\text{student}}(\theta_{t+1}) - \Psi_{\text{student}}(\theta_t)]}_{\text{(B)}}.\]

\textbf{Term (A) --- Gatekeeper update on reference set.} The gatekeeper
update \(S_{t+1} = \Pi_{[0,1]}[S_t + \beta_t(\text{SCXUpdate} - S_t)]\)
moves \(S_t\) toward \(\hat{C}\). However, this movement is \textbf{not
guaranteed} to reduce \(\Psi_{\text{gate}}\) on \(M_0\) because
SCXUpdate uses data from \(M_{t+1}\) (acceptance-biased), not from
\(M_0\). \textbf{Gap: The alignment between SCXUpdate's direction on
\(M_{t+1}\) and the gradient of \(\Psi_{\text{gate}}\) on \(M_0\) is not
established.} Under exploration (C8), \(M_{t+1}\) eventually contains
representatives from all regions, so for large \(t\), the empirical
distributions \(P_{M_{t+1}}\) and \(P_{M_0}\) become close. But the rate
of this convergence depends on the exploration fraction \(\varepsilon\)
and the data distribution's support.

\textbf{Term (B) --- Student update on reference set.} The student is
trained by SGD on \(M_t\) (acceptance-biased), not on \(V_0\). A
standard SGD step reduces the loss on the training distribution
\(P_{S_t}\), but may \textbf{increase} the loss on the reference
distribution \(P_{V_0}\). \textbf{Gap: The relationship between
\(\mathbb{E}_{P_{S_t}}[\ell]\) and \(\mathbb{E}_{P_{V_0}}[\ell]\) is not
controlled.} Under two-timescale separation (C6'), the student takes
many steps between gatekeeper updates, converging approximately to a
local minimum of the current \(P_{S_t}\). The quality of this minimum on
\(P_{V_0}\) depends on the KL divergence between the distributions:
\[\mathbb{E}_{P_{V_0}}[\ell(f_\theta)] \leq \mathbb{E}_{P_{S_t}}[\ell(f_\theta)] + D_{KL}(P_{V_0} \| P_{S_t})/C,\]
which is bounded only if the gatekeeper does not systematically exclude
regions of \(V_0\).

\textbf{Overall assessment.} The descent inequality
\(\mathbb{E}[\Psi_{t+1} \mid \mathcal{F}_t] \leq \Psi_t - \eta_t\) has
\textbf{not been rigorously proven}. The decomposition above identifies
two critical gaps: (i) alignment of gatekeeper update direction on
biased vs.~reference data, and (ii) generalization of student
improvement from training distribution to reference distribution. The
additional conditions (C6', C8, C9) are designed to close these gaps,
but their formal sufficiency has not been established. \textbf{The
descent property is therefore a conjecture, not a theorem.}

\subsubsection{7.5 Downgraded Convergence
Claim}\label{downgraded-convergence-claim}

\textbf{Conjecture SE-1 (formerly ``Theorem SE-1'').} With the explicit
Lyapunov candidate \(\Psi\) defined above and under conditions
(C1')-(C9) (see Document 06 for the full condition set), the SCX
self-evolution sequence satisfies:

\[\Psi(S_t, \theta_t) \xrightarrow{a.s.} \Psi_\infty \geq 0,\]
\[\sum_{t=1}^\infty \bigl( \alpha_t \|\nabla_\theta L_t(\theta_t)\|^2 + \beta_t \|\Delta S_t\|^2 \bigr) < \infty \quad \text{a.s.},\]

and any limit point \((S_\infty, \theta_\infty)\) satisfies the
fixed-point equations. The convergence is in the \textbf{Cesàro mean}
sense rather than monotonic: the Lyapunov function may increase at
individual steps but its long-run average decreases.

\begin{quote}
\textbf{Honest status (2026-06-28).} This is labeled a
\textbf{Conjecture} because the Lyapunov descent property (Definition
24) has not been rigorously proven. The key difficulties are the
selection bias cycle (DEFECT-13) and the tension between gatekeeper
exploration and student stabilization (DEFECT-14). The conditions (C6'),
(C8), and (C9) are proposed mechanisms to address these difficulties,
but their formal sufficiency awaits proof. Numerical evidence is
pending.
\end{quote}

\begin{center}\rule{0.5\linewidth}{0.5pt}\end{center}

\subsection{8. Monotonicity Analysis}\label{monotonicity-analysis}

\subsubsection{8.1 Expected Monotonicity}\label{expected-monotonicity}

\textbf{Corollary 1 (Expected Monotonicity).} Under the conditions of
Theorem 6, even if the losses are not convex but the gradient updates
are unbiased estimators of the true gradient:

\[\mathbb{E}[\Delta V(z_t)] \leq 0 \quad \forall t\]

where the expectation is over the random sampling of batches.

\emph{Proof.} The gradient descent update ensures descent in expectation
for sufficiently small step sizes (standard result for SGD with unbiased
gradients). \(\square\)

\subsubsection{8.2 Strict Monotonicity
Gap}\label{strict-monotonicity-gap}

\textbf{Proposition 11 (Strict Decrease Condition).} The inequality
\(\Delta V(z_t) < 0\) is strict unless both of the following hold: 1.
\(S_t\) is already optimal given \(M_{t+1}\):
\(\nabla_S \hat{L}_{\text{gate}}(S_t; M_{t+1}) = 0\). 2. \(\theta_t\) is
already optimal given \(M_{t+1}\):
\(\nabla_\theta \hat{L}_{\text{nep}}(f_{\theta_t}; M_{t+1}) = 0\).

If either gradient is non-zero, the gradient step strictly decreases the
respective loss (for sufficiently small step sizes), giving
\(\Delta V(z_t) < 0\).

\subsubsection{8.3 Monotonicity of Memory
Bank}\label{monotonicity-of-memory-bank}

A separate monotonicity holds for the memory bank:

\textbf{Proposition 12 (Memory Monotonicity).} The sequence
\(\{N_t\}_{t=0}^\infty\) of memory bank sizes is non-decreasing:

\[N_0 \leq N_1 \leq N_2 \leq \dots\]

and converges to \(N_\infty \leq \infty\).

\emph{Proof.} \(M_t \subseteq M_{t+1}\) (Assumption B1) directly implies
\(N_t \leq N_{t+1}\). Monotone bounded sequences converge (possibly to
\(\infty\)). \(\square\)

\subsubsection{\texorpdfstring{8.4 Relationship Between \(V\) and
\(N_t\)}{8.4 Relationship Between V and N\_t}}\label{relationship-between-v-and-n_t}

\textbf{Proposition 13 (Consistency of Monotonicities).} A decreasing
Lyapunov function \(V(z_t)\) and an increasing memory size \(N_t\) are
not contradictory: \(V\) measures average loss (which decreases as the
system learns), while \(N_t\) measures total accumulated data (which
increases). The system can simultaneously get better (lower \(V\)) and
accumulate more data (higher \(N_t\)).

\begin{center}\rule{0.5\linewidth}{0.5pt}\end{center}

\subsection{9. Phase Portrait}\label{phase-portrait}

\subsubsection{9.1 Qualitative
Description}\label{qualitative-description}

Based on the analysis above, the SCX self-evolution system has the
following qualitative phase portrait:

\textbf{Regime I: Initialization} (\(t = 0\)). - \(S_0\) derived from
static SCX consensus score - \(M_0\) contains initial reliable samples
(or empty) - \(\theta_0\) from pretrained model (or random) - \(V(z_0)\)
is finite and bounded by Theorem 1's guarantee

\textbf{Regime II: Rapid Improvement} (\(0 < t < T_1\)). - The
gatekeeper rapidly improves as it receives feedback from the growing
memory bank - \(V(z_t)\) decreases quickly - The NEP student begins to
provide useful delayed feedback - \(N_t\) grows rapidly

\textbf{Regime III: Diminishing Returns} (\(T_1 \leq t < T_2\)). - The
gatekeeper approaches its optimal performance given the available data -
\(V(z_t)\) decreases slowly - Memory bank growth slows (only hard or
ambiguous samples remain) - The system enters a ``refinement'' phase

\textbf{Regime IV: Saturation} (\(t \geq T_2\)). - \(z_t\) approaches a
fixed point \(z^*\) - \(V(z_t)\) plateaus - \(N_t\) saturates (or grows
only from genuinely novel samples) - The system has converged

\subsubsection{9.2 Parameter Dependence}\label{parameter-dependence}

The qualitative behavior depends on several key parameters:

\begin{longtable}[]{@{}
  >{\raggedright\arraybackslash}p{(\linewidth - 2\tabcolsep) * \real{0.3667}}
  >{\raggedright\arraybackslash}p{(\linewidth - 2\tabcolsep) * \real{0.6333}}@{}}
\toprule\noalign{}
\begin{minipage}[b]{\linewidth}\raggedright
Parameter
\end{minipage} & \begin{minipage}[b]{\linewidth}\raggedright
Effect on Dynamics
\end{minipage} \\
\midrule\noalign{}
\endhead
\bottomrule\noalign{}
\endlastfoot
\(\eta\) (noise rate) & Higher \(\eta\) slows convergence (more noise to
filter) \\
\(\mu_s\) (clean error) & Higher \(\mu_s\) reduces initial \(S_0\)
quality, slowing early improvement \\
\(M\) (number of experts) & Larger \(M\) improves \(S_0\) via tighter
concentration (Theorem 1) \\
\(\delta_{\text{store}}\) & Higher threshold: slower memory growth,
higher precision \\
\(\lambda\) (Lyapunov weight) & Higher \(\lambda\) prioritizes NEP
improvement over gatekeeper \\
\(\alpha, \beta\) (step sizes) & Larger steps: faster convergence but
risk of overshoot \\
\end{longtable}

\subsubsection{9.3 Phase Diagram}\label{phase-diagram}

We can partition the parameter space into regimes:

\textbf{Region A (Convergent)} \(\subset \mathcal{Z}\): Initial states
from which the system converges to a fixed point. Under Assumptions
B1-B3, this region contains all physically realizable initial states.

\textbf{Region B (Limit Cycle)} \(\subset \mathcal{Z}\): Initial states
leading to periodic orbits. \textbf{Conjectured empty} for SCX
(Conjecture 2).

\textbf{Region C (Divergent)} \(\subset \mathcal{Z}\): Initial states
from which the system diverges (e.g., \(N_t \to \infty\) without
convergence of \(S_t\)). This can occur if \(\delta_{\text{store}}\) is
too low, causing the system to accumulate noisy samples indefinitely.

\textbf{Proposition 14 (Divergence Condition).} If
\(\delta_{\text{store}} < \eta\) (the storage threshold is below the
noise rate), the system may diverge in Region C: the gatekeeper stores a
significant proportion of noisy samples, degrading the NEP student,
which further degrades the gatekeeper, creating a vicious cycle.

\emph{Proof sketch.} If \(\delta_{\text{store}} < \eta\), then even a
random gatekeeper would store \textgreater{} \(\eta\) fraction of noisy
samples. The stored noise contaminates \(M_{t+1}\), which trains
\(\theta_{t+1}\), potentially amplifying the error. Formalizing this
requires analyzing the noise propagation eigenvalue (see Proposition
15). \(\square\)

\begin{center}\rule{0.5\linewidth}{0.5pt}\end{center}

\subsection{10. Connection to Existing Theorems
1-3}\label{connection-to-existing-theorems-1-3}

\subsubsection{10.1 Theorem 1: Noise Detection
Guarantee}\label{theorem-1-noise-detection-guarantee}

Theorem 1 provides the \textbf{initialization guarantee} for the
dynamical system:

\[V(z_0) \leq 1 - \text{F1}_0 \leq \frac{1}{\eta} \sum_{s \in \mathcal{S}} \rho_s \cdot \exp\bigl(-2M\Delta_s^2\bigr)\]

This bounds the initial Lyapunov value \(V(z_0)\) in terms of the noise
detection F1 score.

\textbf{Dynamical implication}: The system starts with
\(V(z_0) \leq \varepsilon_0(M)\) where \(\varepsilon_0(M)\) decreases
exponentially in \(M\). A larger \(M\) gives a better initialization,
reducing the number of evolution rounds needed to reach a given
performance level.

\subsubsection{10.2 Theorem 2: Weak Feature Failure Lower
Bound}\label{theorem-2-weak-feature-failure-lower-bound}

Theorem 2 constrains the \textbf{improvement rate} of the dynamical
system:

If \(\phi\) is \(\delta\)-weak, then for any \(t\):

\[F1_t \leq F1_{\text{base}} + C_F \sqrt{\frac{\delta}{2}} + \text{(evolution improvement)}\]

The ``evolution improvement'' term captures what self-evolution adds
beyond the static baseline. \textbf{Conjectured behavior}: this
improvement term itself is bounded by \(O(\sqrt{\delta})\) if the state
space is discovered from \(\phi\), because the NEP student cannot
extract state information that isn't in the features.

\textbf{Corollary 2 (Improvement Bound Under Weak Features).} If the
feature mapping \(\phi\) is \(\delta\)-weak and the state discovery uses
\(\phi\), then the improvement from self-evolution is bounded by:

\[\limsup_{t \to \infty} (F1_t - F1_0) \leq C \cdot \sqrt{\frac{\delta}{2}}\]

for some constant \(C > 0\) depending on \(\eta\) and the Lipschitz
constant of the F1 score.

\emph{Status: \textbf{Conjecture} --- a rigorous proof requires showing
that self-evolution cannot amplify information not present in the
features. The information bottleneck principle suggests this is true.}

\subsubsection{10.3 Theorem 3: Noise-Difficulty
Unidentifiability}\label{theorem-3-noise-difficulty-unidentifiability}

Theorem 3 establishes a fundamental \textbf{obstacle} to convergence to
the perfect fixed point:

If the data generating process satisfies the unidentifiability
conditions, then even the limiting gatekeeper \(S_\infty\) cannot
achieve perfect distinction between noise and difficulty.

\textbf{Proposition 15 (Residual Error at Fixed Point).} Under the
conditions of Theorem 3, any fixed point \(z^*\) of the SCX
self-evolution system must satisfy:

\[\mathbb{E}_{x \in s_1}\bigl[ |S^*(x, y_{\text{obs}}) - \mathbf{1}\{y_{\text{obs}} = f^*(x)\}| \bigr] \geq \frac{\eta\rho}{2}\]

where \(s_1\) is the ambiguous state from the Theorem 3 construction.

\emph{Proof.} Theorem 3 shows that in the ambiguous subset, any
algorithm has error rate \(\geq \eta\rho/2\). Since \(S^*\) induces a
noise detection algorithm via thresholding, the same lower bound
applies. \(\square\)

\textbf{Corollary 3 (Non-Vanishing Lyapunov Function).} Under the
conditions of Theorem 3, the Lyapunov function cannot converge to zero:

\[\liminf_{t \to \infty} V(z_t) \geq c(\eta, \rho) > 0\]

where
\(c(\eta, \rho) = \eta \cdot \rho \cdot \bigl( \log(1/\eta) + \log(2/\rho) \bigr) / 2\)
(approximately).

\emph{Status: \textbf{Conjecture} --- the exact constant
\(c(\eta, \rho)\) depends on the loss function and the data
distribution. The existence of a positive lower bound follows from
Proposition 15, but the precise value is not yet derived.}

\subsubsection{10.4 Integration}\label{integration}

The three theorems together paint a complete picture of the dynamical
system:

\begin{enumerate}
\def\labelenumi{\arabic{enumi}.}
\tightlist
\item
  \textbf{Theorem 1} (initialization): The system starts close to the
  optimal fixed point when \(M\) is large.
\item
  \textbf{Theorem 2} (speed limit): Weak features constrain how much
  self-evolution can improve over the initialization.
\item
  \textbf{Theorem 3} (fundamental barrier): Even with infinite
  evolution, perfect noise detection is impossible.
\end{enumerate}

The self-evolution dynamics bridge these three fixed points:
initialization → improvement (bounded by feature quality) → saturation
(bounded by unidentifiability).

\begin{center}\rule{0.5\linewidth}{0.5pt}\end{center}

\subsection{11. Summary of Proven vs.~Conjectured
Claims}\label{summary-of-proven-vs.-conjectured-claims}

\begin{longtable}[]{@{}
  >{\raggedright\arraybackslash}p{(\linewidth - 4\tabcolsep) * \real{0.2800}}
  >{\raggedright\arraybackslash}p{(\linewidth - 4\tabcolsep) * \real{0.3200}}
  >{\raggedright\arraybackslash}p{(\linewidth - 4\tabcolsep) * \real{0.4000}}@{}}
\toprule\noalign{}
\begin{minipage}[b]{\linewidth}\raggedright
Claim
\end{minipage} & \begin{minipage}[b]{\linewidth}\raggedright
Status
\end{minipage} & \begin{minipage}[b]{\linewidth}\raggedright
Evidence
\end{minipage} \\
\midrule\noalign{}
\endhead
\bottomrule\noalign{}
\endlastfoot
Proposition 3 (causal structure) & \textbf{Proven} & By inspection of
update definitions \\
Proposition 4 (completeness) & \textbf{Proven} & Standard topology
results \\
Proposition 5 (compactness) & \textbf{Proven} & Arzela-Ascoli under
Lipschitz bound \\
Proposition 6 (bounded orbits) & \textbf{Proven} & From B1-B3 \\
Proposition 7 (fixed points as attractors) & \textbf{Proven} & By
Lyapunov's theorem \\
Proposition 8 (trivial fixed points) & \textbf{Proven} & By
verification \\
Theorem 4 (general existence) & \textbf{Proven} & Brouwer/Schauder fixed
point theorem \\
Theorem 5 (unique existence, convex) & \textbf{Proven} & Contraction
mapping / convex analysis \\
Proposition 9 (necessary condition) & \textbf{Proven} & By
contradiction \\
Proposition 10 (Lyapunov properties) & \textbf{Proven} & Non-negativity,
boundedness, continuity of the explicit \(\Psi\) are verified for finite
\(M_0\) \\
Theorem 6 (Lyapunov decrease) & \textbf{CONJECTURED (was ``Proven'')} &
See DEFECT-03/04. The original proof assumed memory bank growth always
improves losses and gradient steps always descend --- this ignores the
selection bias cycle and distribution shift. The corrected analysis
identifies two gaps: (i) alignment of gatekeeper update direction on
biased vs.~reference data, and (ii) generalization of student
improvement across distributions. Descent is conjectured under
additional conditions (C6', C8, C9). \\
Corollary 1 (expected monotonicity) & \textbf{CONJECTURED} (depends on
Theorem 6) & Conditional on the descent property being proven \\
Proposition 11 (strict decrease) & \textbf{CONJECTURED} (depends on
Theorem 6) & Conditional on Theorem 6 and stationarity conditions \\
Proposition 12 (memory monotonicity) & \textbf{Proven} & From B1 \\
Proposition 13 (consistency) & \textbf{Proven} & By definitions \\
Proposition 14 (divergence condition) & \textbf{Rigorous sketch} &
Requires formalizing noise propagation \\
Proposition 15 (residual error) & \textbf{Proven} & From Theorem 3 \\
\textbf{Conjecture 2} (no periodic orbits) & \textbf{Conjecture} &
Requires extending Lyapunov argument to non-convex \\
\textbf{Conjecture 3} (unique attractor) & \textbf{Conjecture} &
Function-space uniqueness not guaranteed \\
\textbf{Corollary 2} (improvement bound) & \textbf{Conjecture} &
Information bottleneck extension of Theorem 2 \\
\textbf{Corollary 3} (non-vanishing V) & \textbf{Conjecture} & Exact
constant not yet derived \\
\end{longtable}

\begin{center}\rule{0.5\linewidth}{0.5pt}\end{center}

\subsection{References}\label{references}

The dynamical systems concepts used here are standard. Key references
for the mathematical framework:

\begin{enumerate}
\def\labelenumi{\arabic{enumi}.}
\tightlist
\item
  Strogatz, S. H. (2018). \emph{Nonlinear Dynamics and Chaos} (2nd ed.).
  CRC Press. --- Phase portrait analysis, Lyapunov functions.
\item
  Hale, J. K. (1988). \emph{Asymptotic Behavior of Dissipative Systems}.
  AMS. --- \(\omega\)-limit sets, attractors.
\item
  Khalil, H. K. (2002). \emph{Nonlinear Systems} (3rd ed.). Prentice
  Hall. --- Lyapunov stability, invariance principles.
\item
  Conley, C. (1978). \emph{Isolated Invariant Sets and the Morse Index}.
  CBMS Regional Conference Series. --- Lyapunov functions for discrete
  systems.
\item
  Benaïm, M. (1999). ``Dynamics of stochastic approximation
  algorithms.'' \emph{Le Seminaire de Probabilites}, 1709, 1-68. ---
  Stochastic approximation and ODE methods for discrete dynamics.
\item
  Polyak, B. T., \& Juditsky, A. B. (1992). ``Acceleration of stochastic
  approximation by averaging.'' \emph{SIAM Journal on Control and
  Optimization}, 30(4), 838-855. --- Convergence rates for iterative
  methods.
\end{enumerate}

\begin{center}\rule{0.5\linewidth}{0.5pt}\end{center}

\emph{End of Document 02: SCX Self-Evolution as a Discrete Dynamical
System}

\emph{Next: Document 03 analyzes the online learning regret of the
gatekeeper under delayed feedback from the NEP student.}

\begin{center}\rule{0.5\linewidth}{0.5pt}\end{center}

\subsection{Changelog}\label{changelog}

\begin{longtable}[]{@{}
  >{\raggedright\arraybackslash}p{(\linewidth - 6\tabcolsep) * \real{0.1875}}
  >{\raggedright\arraybackslash}p{(\linewidth - 6\tabcolsep) * \real{0.2500}}
  >{\raggedright\arraybackslash}p{(\linewidth - 6\tabcolsep) * \real{0.2500}}
  >{\raggedright\arraybackslash}p{(\linewidth - 6\tabcolsep) * \real{0.3125}}@{}}
\toprule\noalign{}
\begin{minipage}[b]{\linewidth}\raggedright
Date
\end{minipage} & \begin{minipage}[b]{\linewidth}\raggedright
Defect
\end{minipage} & \begin{minipage}[b]{\linewidth}\raggedright
Change
\end{minipage} & \begin{minipage}[b]{\linewidth}\raggedright
Severity
\end{minipage} \\
\midrule\noalign{}
\endhead
\bottomrule\noalign{}
\endlastfoot
2026-06-28 & DEFECT-03 & \textbf{Explicit Lyapunov function defined.}
Replaced the vague composite loss
\(V(z_t) = \mathbb{E}[L_{\text{gate}}] + \lambda \mathbb{E}[L_{\text{nep}}]\)
with a concrete candidate
\(\Psi(S_t, \theta_t) = \frac{1}{|M_0|}\sum_{x \in M_0}(S_t - \hat{C})^2 + \frac{\lambda}{|V_0|}\sum_{V_0}\ell(f_{\theta_t}, y)\)
evaluated on a fixed reference set \(M_0\). The function is
non-negative, bounded, and well-defined for finite \(M_0\). & FATAL \\
2026-06-28 & DEFECT-04 & \textbf{Honest CONJECTURE status for descent
property.} The Lyapunov descent proof (Theorem 6) is acknowledged as
\textbf{not proven}. Two critical gaps identified: (i) alignment of
gatekeeper update direction on biased vs.~reference data, and (ii)
generalization of student improvement across distributions. Theorem 6 is
now labeled CONDITIONAL, CONJECTURED. Corollary 1 and Proposition 11
downgraded to CONJECTURED (depend on Theorem 6). & FATAL \\
2026-06-28 & DEFECT-03/13 & \textbf{Selection bias confound
acknowledged.} Section 7.4 explains why the original Lyapunov candidate
(evaluated on acceptance-biased distribution) could decrease without
genuine improvement. The reference-set fix (evaluating on \(M_0\),
\(V_0\)) eliminates this confound. Additional conditions (C6':
two-timescale, C8: exploration, C9: annealing threshold) proposed as
mitigation. & FATAL/MAJOR \\
2026-06-28 & DEFECT-13/14 & \textbf{Cross-reference to selection bias
analysis and two-timescale scheduling} in Document 06
(06\_fixed\_point\_convergence.md Sections 5.1 and 3.3). & MAJOR \\
\end{longtable}

\end{document}
